% SPDX-License-Identifier: GPL-3.0-or-later
% Copyright (C) 2026 Alessandro Greco
\chapter{Censorship, Bridges, and Pluggable Transports}\label{cap:16-censura}
\aggiornatoal{23 September 2026}{tor 0.4.9.13; pt-spec v1; obfs4proxy 0.0.14 (Debian); Tor Project documentation}

\begin{obiettivi}
By the end of this chapter, you will be able to:
\begin{itemize}[nosep]
  \item describe the ways a censor can block Tor;
  \item explain what bridges are and why distributing them is difficult;
  \item describe the protocol through which tor communicates with a pluggable transport;
  \item distinguish the strategies used by obfs4, meek, Snowflake, WebTunnel, and Conjure;
  \item operate an obfs4 bridge on the laboratory network and observe what a censor can see (Laboratory~10).
\end{itemize}
\end{obiettivi}

Everything we have covered so far assumes that the client can reach the network. In many countries it cannot: connections to Tor are blocked, sometimes only during periods of political tension and sometimes permanently. This chapter concludes Part~II with \textbf{censorship circumvention}, one of the Tor Project's most active areas in recent years.

% ---------------------------------------------------------------------------
\section{How Tor Is Blocked}\label{sec:cens-blocco}

A censor, usually a government controlling a country's connectivity providers, has several tools:
\begin{description}
  \item[Address blocking] The relay list is public: it is the consensus (\cref{cap:12-directory}). A censor can simply download it and block every relay IP address, or just the directory authorities and relays from which clients obtain consensuses.
  \item[Website blocking] Blocking \texttt{torproject.org} prevents people from downloading Tor Browser, even if the Tor network itself remains reachable.
  \item[Protocol recognition] \emph{Deep packet inspection} (DPI) may recognize a Tor connection from properties of its TLS handshake, packet sizes, or timing, even when the destination address is not yet known.
  \item[Active probing] When DPI suggests that a server may be an undisclosed relay, the censor can connect to it and try speaking the Tor protocol. If the server responds, it can be blocked.
  \item[Allowlisting] In extreme cases, a censor permits only explicitly authorized protocols or destinations and blocks everything else.
\end{description}
Each technique has a corresponding countermeasure, and the contest between censors and developers continues. Projects such as OONI measure where and how Tor is blocked (\cref{cap:24-metrics}).

% ---------------------------------------------------------------------------
\section{Bridges}\label{sec:cens-bridge}

The response to address blocking is the \textbf{bridge}: a relay absent from the public consensus. A bridge acts as a guard (\cref{sec:guard}). The client uses it as the first hop, and the circuit then continues normally through the network. Because the bridge address is not public, a censor cannot block it merely by reading a list. Bridges publish descriptors to a \textbf{bridge authority}, which checks their reachability and maintains a list without making it public. Published bridge statistics use a digest \emph{of a digest} of the identity to avoid exposing it.

A client configures a bridge using a \emph{bridge line}. On the laboratory network, client \texttt{test011bc} has the following in its \file{torrc}:
\begin{lstlisting}[language=torrc,numbers=none]
UseBridges 1
Bridge obfs4 127.0.0.1:9909 2557828B596B33AB02075FE7EAD86F2DC3CEAFD5 cert=4sKDQih... iat-mode=0
ClientTransportPlugin obfs4 exec /usr/bin/obfs4proxy
\end{lstlisting}
The line specifies the transport type, address, bridge identity fingerprint (used to verify the peer, \cref{cap:07-canali}), and transport parameters. The value \texttt{cert=4sKDQih...} is \emph{abbreviated} for print and must not be copied as a valid configuration. Readers must use the complete certificate and fingerprint generated by their own bridge instance.

\subsection{The distribution problem}

A bridge is useful for as long as the censor does not know its address. But to be useful, its address must be given to users, and a censor can pretend to be a user. There is no perfect solution: a distributor must give enough bridges to many legitimate users while revealing as few as possible to the censor. The Tor Project distributes bridges through several channels, each with its own defenses against automated harvesting~\cite{torsupport}:
\begin{itemize}
  \item the website \url{https://bridges.torproject.org/};
  \item email, by writing to \texttt{bridges@torproject.org};
  \item a Telegram bot (\texttt{@GetBridgesBot});
  \item directly through Tor Browser, which detects blocking and requests suitable bridges for the user's country (\emph{Connection Assist}, based on the channel called \emph{moat});
  \item private channels run by organizations that know their users personally.
\end{itemize}
For years, the system assigning bridges to distribution channels was BridgeDB. It was replaced by \textbf{rdsys}, with migration completed in October 2024~\cite{torblogrdsys}. Bridge operators can select a channel with the \opt{BridgeDistribution} option.

\subsection{Lox}

\textbf{Lox}~\cite{tulloch2023lox} approaches distribution with a \emph{reputation} system. A new user receives a bridge and an anonymous credential. If that bridge remains unblocked over time, the credential gains levels, up to a maximum of 4; starting at level 2, it also permits the user to invite others. If the bridge is blocked, the user can obtain another but loses reputation. The premise is that a censor that blocks every bridge it receives cannot gain levels, whereas honest users can. The credentials employ cryptography that prevents even the system itself from learning who invited whom, protecting users' \emph{social graph}. Lox is being tested in alpha versions of Tor Browser.

% ---------------------------------------------------------------------------
\section{Pluggable Transports}\label{sec:cens-pt}

Hiding the address is not enough to defeat protocol recognition: the traffic itself must look different. Rather than modifying tor whenever a censor learns to recognize a new feature, the Tor Project has defined an interface for external programs that transform traffic: \textbf{pluggable transports} (PTs)~\spec{pt-spec}{Architecture Overview}.

\subsection{The architecture}

A PT is a separate process launched by tor (hence \emph{managed}). On the client side, it offers tor a local SOCKS proxy. tor hands the connection to the bridge to the PT, which transforms the traffic before sending it. On the bridge side, a PT listens on the public port, reverses the transformation, and forwards traffic to tor's local ORPort. It may use a small additional protocol, the \emph{Extended ORPort}, to tell tor the client's actual address for statistics.

\begin{figure}[htbp]
\centering
% SPDX-License-Identifier: GPL-3.0-or-later
% Copyright (C) 2026 Alessandro Greco
%% Accessibility alternative text, used only in EPUB.
\altfig{Simplified pluggable transport architecture: tor connects to the client PT's local SOCKS proxy, which transforms traffic. At the bridge, the server PT reverses the transformation and forwards it to tor via ORPort or, when configured, Extended ORPort. Snowflake uses volunteer intermediary proxies not drawn here. A censor observes transformed traffic on the monitored link, but might still detect it. Examples: obfs4 resembles random bytes, Snowflake uses WebRTC, and WebTunnel resembles HTTPS.}
% Pluggable transport architecture (pt-spec)
\begin{tikzpicture}[font=\sffamily\scriptsize,
  proc/.style={draw,rounded corners=2pt,align=center,minimum height=0.8cm,minimum width=1.7cm},
  mac/.style={draw=torgrigio,dashed,rounded corners=4pt,inner sep=6pt}]
  \node[proc,fill=torviola!12] (tc) at (0,0) {tor\\(client)};
  \node[proc,fill=torverde!15] (pc) at (2.6,0) {client PT};
  \node[proc,fill=torverde!15] (ps) at (8.2,0) {server PT};
  \node[proc,fill=torviola!12] (tb) at (10.8,0) {tor\\(bridge)};
  \node[mac,fit=(tc)(pc),label={[font=\sffamily\scriptsize]above:user's computer}] {};
  \node[mac,fit=(ps)(tb),label={[font=\sffamily\scriptsize]above:bridge}] {};
  \draw[->] (tc) -- node[above,font=\sffamily\tiny] {SOCKS} (pc);
  \draw[->] (ps) -- node[above,align=center,font=\sffamily\tiny] {ORPort /\\Extended ORPort} (tb);
  \draw[->,thick,decorate,decoration={snake,amplitude=1.2pt,segment length=6pt}] (pc) -- node[above] {transformed traffic} (ps);
  \node[server,minimum width=1.4cm,fill=avviso!10,draw=avviso] (cen) at (5.4,-1.5) {censor};
  \draw[dotted,avviso] (cen) -- (5.4,-0.15);
  \node[text=torgrigio,align=center] at (5.4,-2.3) {random-looking bytes (obfs4),\\WebRTC (Snowflake), HTTPS (WebTunnel)};
  \node[text=torgrigio,font=\sffamily\tiny,align=center,text width=10.5cm] at (5.4,-3.05) {Simplified: Snowflake uses intermediary volunteer proxies not shown};
\end{tikzpicture}

\caption{Schematic pluggable transport architecture. On the client, tor communicates with the PT via a local SOCKS proxy; at the bridge, the PT forwards restored traffic to tor via ORPort or, when configured, Extended ORPort. Snowflake also involves volunteer intermediary proxies, omitted here. A censor sees transformed traffic on the monitored link, but it may still be detectable.}
\label{fig:pt}
\end{figure}

\subsection{The protocol between tor and the PT}

tor configures the PT through environment variables, and the PT responds by writing lines to its standard output~\spec{pt-spec}{Specification}. We can test this by manually starting \cmd{obfs4proxy} in the laboratory container, just as tor would. As a client:
\begin{lstlisting}[language=bash,numbers=none]
$ env TOR_PT_MANAGED_TRANSPORT_VER=1 TOR_PT_STATE_LOCATION=/tmp/pt \
      TOR_PT_EXIT_ON_STDIN_CLOSE=1 TOR_PT_CLIENT_TRANSPORTS=obfs4 obfs4proxy
VERSION 1
CMETHOD obfs4 socks5 127.0.0.1:40057
CMETHODS DONE
\end{lstlisting}
The PT declares the protocol version and announces that it provides the \texttt{obfs4} transport as a SOCKS5 proxy on a local port of its own choosing. As a server, with variables specifying the public listening address and the ORPort destination:
\begin{lstlisting}[language=bash,numbers=none]
$ env ... TOR_PT_SERVER_TRANSPORTS=obfs4 \
      TOR_PT_SERVER_BINDADDR=obfs4-127.0.0.1:19999 \
      TOR_PT_ORPORT=127.0.0.1:5104 obfs4proxy
VERSION 1
SMETHOD obfs4 127.0.0.1:19999 ARGS:cert=5W6vcbWbFGGZgabwm4C3WipezwnAfadoBhqfNVpY72xnkjB/Szczbxg8fLWJJ1uRlyjPSg,iat-mode=0
SMETHODS DONE
\end{lstlisting}
The parameters following \texttt{ARGS:} are what the client will need; they appear in the bridge line. The PT also writes them to \file{obfs4\_bridgeline.txt} in its state directory. The \file{TOR\_PT\_EXIT\_ON\_STDIN\_CLOSE} variable tells the PT to terminate when tor closes its input, so the PT cannot outlive the process that launched it. (Through an oversight, the specification's example says the PT writes its messages to \enquote{stdin}; in practice, as we have just seen, it writes them to standard output.)

The server command's \texttt{env ...} is \emph{schematic}: the ellipsis stands for other required environment variables and must not be typed literally. Likewise, the SOCKS5 port in the client output comes from a reference session and is not fixed.

\subsection{Lyrebird}

The transport program used by Tor Browser is \textbf{lyrebird}, derived from the earlier \cmd{obfs4proxy}. It implements obfs4, meek, Snowflake, and WebTunnel~\cite{torsupport}. The laboratory uses \cmd{obfs4proxy} 0.0.14, available in Debian packages and still compatible with obfs4.

% ---------------------------------------------------------------------------
\section{Strategies}\label{sec:cens-strategie}

Transports currently in use follow two opposing strategies: \emph{look like nothing recognizable} or \emph{look like something permitted}.

\subsection{obfs4: looking like nothing recognizable}

\textbf{obfs4}, designed by Yawning Angel, transforms traffic into a sequence of bytes indistinguishable from random data: no recognizable header, no TLS handshake, and variable packet sizes. Three design choices make it effective:
\begin{itemize}
  \item Its key exchange is a variant of ntor (\cref{sec:ntor}) in which Curve25519 public keys are encoded with Elligator2, a technique that makes a curve point indistinguishable from a random string.
  \item The client must know the bridge's \texttt{cert} parameter, which contains its identity and public key. Without that parameter, an observer cannot construct a valid handshake and the bridge does not respond. This defeats \textbf{active probing}: a curious censor connecting to the bridge receives no recognizable answer.
  \item The \texttt{iat-mode} (\emph{inter-arrival time}) parameter can, at a performance cost, add randomized delays and fragmentation to obscure timing as well.
\end{itemize}
The limitation of obfs4 is also its strength: a censor employing an allowlist, or blocking \enquote{anything that looks random}, can stop it. This happens in some countries.

\subsection{meek: behind a large service}

\textbf{meek}~\cite{fifield2015fronting} uses \emph{domain fronting}. The client opens an HTTPS connection to a large content-distribution service, using an innocuous domain name in the visible part of the connection, and specifies the true destination (a server acting as a gateway to Tor) only within encrypted data. Blocking meek would require the censor to block the entire service, with extensive collateral damage. However, this technique depends on providers permitting it, and many have made it impossible in recent years. meek is also slow and expensive because all traffic passes through the provider's servers.

\subsection{Snowflake: many ephemeral proxies}

\textbf{Snowflake}~\cite{bocovich2024snowflake} uses thousands of volunteers running small proxies, for example a browser extension or a program on a server. The client contacts a \emph{broker} through a hard-to-block channel such as domain fronting to obtain a proxy, then connects to it using WebRTC, the protocol used for video calls in browsers. The proxy forwards traffic to a Snowflake bridge run by the Tor Project. There are many proxies, and they change constantly. Blocking all their addresses is impractical, while blocking WebRTC means blocking video calls.

\subsection{WebTunnel: looking like a website}

\textbf{WebTunnel}, announced by the Tor Project in March 2024, wraps traffic in an HTTPS connection with behavior similar to WebSocket~\cite{torblogwebtunnel}. The bridge can sit behind an ordinary website, sharing its domain, address, and port. A reverse proxy forwards only requests for a secret path to the bridge and serves the normal website to everyone else. WebTunnel targets censors using allowlists, against which obfs4 does not work: it has the \enquote{shape} of a permitted protocol.

\subsection{Conjure: hiding in the address space}

\textbf{Conjure}~\cite{frolov2019conjure} is a \emph{refraction networking} system. Instead of a bridge at a known address, it uses stations installed at cooperating connectivity providers. These stations observe traffic destined for ranges of unused addresses and divert it toward Tor. The client connects to a \enquote{phantom} address chosen from millions of possibilities. Blocking Conjure would require the censor to block entire portions of the Internet. According to Tor Project documentation, Conjure has been tested successfully but still needs stabilization work and is available only in alpha versions of Tor Browser and some applications~\cite{torsupport}.

\begin{attenzione}[No transport works everywhere]
Every transport has weaknesses, and effectiveness changes by country and from month to month. In the Tor Project forum, for example, users reported in April 2025 that obfs4, meek, and Snowflake were unusable in China and that WebTunnel connected but was quickly blocked~\cite{forumcina2025}. Tor Browser therefore offers several transports and a feature that suggests one suited to the user's country.
\end{attenzione}

% ---------------------------------------------------------------------------
\section{Laboratory: An obfs4 Bridge Through a Censor's Eyes}\label{sec:lab-censura}

\begin{laboratorio}[10]{An obfs4 bridge on the laboratory network}
\textbf{Objective.} Start a network containing an obfs4 bridge and a client using it; capture and compare traffic from an ordinary Tor connection and an obfs4 connection.\\
\textbf{Prerequisites.} The laboratory network (\cref{cap:19-rete-privata}). The image includes \cmd{obfs4proxy} and \cmd{snowflake-client} from Debian packages.\\
\textbf{Estimated time.} 1--2 hours.
\end{laboratorio}

\paragraph{Step 1: the bridge network.} The \file{labs/en/lab/chutney/networks/book-bridge} file adds to the book's network a bridge authority (\texttt{test003ba}), an obfs4 bridge (\texttt{test009br}), and a client using it (\texttt{test011bc}):
\begin{lstlisting}[language=bash,numbers=none]
$ make -C labs/en/lab clean
$ make -C labs/en/lab start NETWORK=book-bridge
\end{lstlisting}
The client's log from an earlier reference session includes:
\begin{lstlisting}[numbers=none]
[notice] new bridge descriptor 'test009br' (fresh): $2557828B...~test009br
[warn] Only one bridge (transport: 'obfs4') is configured. You should
       have at least two for conflux, for any transport that is not 'snowflake'.
[notice] Bootstrapped 100% (done): Done
\end{lstlisting}
The client obtained the bridge descriptor through the bridge itself and finished bootstrapping. The warning reminds us that Conflux (\cref{cap:31-futuro}) uses two parallel circuits and cannot do so with just one bridge. A web request through this client's SOCKS port (9011) works as it does with the ordinary client.

\paragraph{Step 2: what the censor sees.} Inside the container, we capture traffic with \cmd{tcpdump} on the bridge's obfs4 port (9909) and a normal ORPort (5104), while the client uses the bridge and we open a channel to \texttt{test004r} with Laboratory~1. We then analyze the capture with \cmd{tshark}. Here is the result from one of our sessions:
\begin{lstlisting}[numbers=none]
== ORPort 5104 (TLS), first bytes sent by the client ==
1533  16030105f8010005f403033e1234a0401520d1c52b...
== obfs4 9909, first bytes sent by the client ==
1957  b566475349099523c53005e1004a392160f15d473e...
== packets recognized as TLS ==
port 5104: 83
port 9909: 0
== obfs4 segment sizes (client -> bridge) ==
1957 1677 1888 1891 1716 1716 3436 1395 1716 1748 1500 1194 ...
\end{lstlisting}
On the ORPort, the first bytes are unmistakable: \texttt{16 03 01} starts a TLS handshake record, and Wireshark recognizes the entire conversation as TLS. The censor sees a TLS connection to inspect for Tor-specific features; the relay's address is also listed in the consensus. Toward the obfs4 bridge, the bytes have no recognizable structure, no packet is identified as TLS, and segment sizes do not mirror 514-byte cells. The censor sees only a stream of apparently random data toward an unknown address.

\paragraph{Step 3: probing the bridge.} We try connecting to the bridge's port 9909 using Laboratory~1, as a censor suspecting a relay might:
\begin{lstlisting}[numbers=none]
TimeoutError: The handshake operation timed out
\end{lstlisting}
The bridge does not respond \emph{at all}: no error message and no bytes, until our client gives up after 15 seconds. To obfs4, our TLS ClientHello is an invalid handshake, and the response to an invalid handshake is silence. A censor probing the address learns nothing, not even that a service speaking some other protocol is listening there. Without the \texttt{cert} parameter, it cannot construct a valid obfs4 handshake anyway.

% ---------------------------------------------------------------------------
\section*{Summary}
\begin{itemize}
  \item A censor can block Tor by address, recognize its protocol, probe suspect servers, or permit only explicitly authorized traffic.
  \item Bridges are unpublished relays. Their distribution, currently handled by rdsys and being explored with Lox's reputation mechanism, balances access for users against discovery by censors.
  \item Pluggable transports are external processes that transform traffic; tor configures them through environment variables and reads their responses (\texttt{CMETHOD}, \texttt{SMETHOD}).
  \item obfs4 makes traffic look random and resists active probing; meek, Snowflake, WebTunnel, and Conjure conceal it behind services, protocols, or addresses that are costly for a censor to block.
  \item No transport works everywhere or forever.
\end{itemize}

% ---------------------------------------------------------------------------
\section*{Exercises}

\begin{esercizio}[1]\label{es:16-guard}
How does a bridge differ from a guard? Why does a client using bridges not also use a public guard?
\end{esercizio}

\begin{esercizio}[1]\label{es:16-strategie}
Classify obfs4, meek, Snowflake, WebTunnel, and Conjure by strategy (\enquote{look like nothing recognizable} or \enquote{look like something permitted}) and by what a censor would have to block to stop each one.
\end{esercizio}

\begin{esercizio}[2]\label{es:16-sondaggio}
Describe how a censor might discover an obfs4 bridge by observing only traffic, without actively probing it. Which statistical features could it use, and what does \texttt{iat-mode} do to counter them?
\end{esercizio}

\begin{esercizio}[2]\label{es:16-lox}
In Lox, why can a censor not advance through reputation levels? What could it do to evade the mechanism, and what makes that costly?
\end{esercizio}

\begin{esercizio}[2]\label{es:16-entropia}
Calculate the Shannon entropy (in bits per byte) of the first 10\,KB captured from an obfs4 connection and an ordinary TLS connection in the laboratory. What do you expect, and why might the difference be small? What, then, makes TLS recognizable?
\end{esercizio}

\begin{esercizio}[3]\label{es:16-sondaggio-lab}
Try running Laboratory~1 against the obfs4 port of your bridge. What error occurs? Then use \cmd{obfs4proxy} as a client with an incorrect \texttt{cert} parameter: what do you observe in the bridge logs?
\end{esercizio}

% ---------------------------------------------------------------------------
\section*{Further Reading}
\begin{itemize}
  \item \titolo{Pluggable Transport Specification (Version 1)} (pt-spec)~\cite{torspec}.
  \item L.~Tulloch and I.~Goldberg, \titolo{Lox: Protecting the Social Graph in Bridge Distribution}~\cite{tulloch2023lox}.
  \item D.~Fifield et al., \titolo{Blocking-resistant communication through domain fronting}~\cite{fifield2015fronting}.
  \item C.~Bocovich et al., \titolo{Snowflake, a censorship circumvention system using temporary WebRTC proxies}~\cite{bocovich2024snowflake}.
  \item S.~Frolov et al., \titolo{Conjure: Summoning Proxies from Unused Address Space}~\cite{frolov2019conjure}.
  \item The Tor Project blog post announcing WebTunnel~\cite{torblogwebtunnel} and the bridge support pages~\cite{torsupport}.
\end{itemize}
